% ===== Document Preamble =====
\documentclass[11pt]{article}
\usepackage[margin=0.5in]{geometry}
\usepackage{amsmath,amssymb}
\usepackage{xcolor}
\usepackage{enumitem}
\usepackage{needspace}
\usepackage{comment}

% ===== Colors =====
\definecolor{AnsRed}{RGB}{170,0,0}
\definecolor{WrongBlue}{RGB}{0,0,160}

% ===== Choice list styling (tight, consistent) =====
\setlist[enumerate]{%
  leftmargin=1.5em,
  itemsep=0.25em,
  topsep=0.25em,
  label=\textbf{\Alph*.}
}

% ===== Style Wrappers =====
\newcommand{\QuestionStyle}[1]{{\large\bfseries\boldmath #1}}
\newcommand{\AnswerStyle}[1]{{\large\bfseries\boldmath #1}}
\newcommand{\ReasonStyle}[1]{{\normalsize #1}}

% ===== Show / Hide Reasoning Toggle =====
\newif\ifshowreason
\showreasontrue      % Instructor version
% \showreasonfalse   % Student version

% ===== Choice Item Macros (Answer + Reasoning) =====
% Usage:
%   \Correct{Answer}{Reasoning}
%   \Wrong{Answer}{Reasoning}

\newcommand{\Correct}[2]{%
  \item \AnswerStyle{\textcolor{AnsRed}{#1}}%
  \ifshowreason\\[-0.2em]
    \ReasonStyle{\textcolor{AnsRed}{#2}}%
  \fi
}

\newcommand{\Wrong}[2]{%
  \item \AnswerStyle{\textcolor{WrongBlue}{#1}}%
  \ifshowreason\\[-0.2em]
    \ReasonStyle{\textcolor{WrongBlue}{#2}}%
  \fi
}

% ===== Question Block Macro =====
% Usage: \QBlock{Q1}{stem...}{<items>}

\newcommand{\QBlock}[3]{%
  \Needspace{8\baselineskip}%
  \noindent
  \begin{minipage}{\linewidth}
    \QuestionStyle{#1.~#2}

    \vspace{0.25em}
    \begin{enumerate}
      #3
    \end{enumerate}

    \vspace{0.25em}
    \hrule height 0.4pt
  \end{minipage}
  \vspace{0.75em}
}

% ===== Document Start =====

\begin{document}

\begin{center}
{\huge \textbf{Introductory Physics Isomorphic Questions}}\\[0.25em]
{\Large Modeling, Units, Vectors, Kinematics, and Dynamics}\\[0.25em]
{\large Multiple choice questions generated below. Answers in RED and incorrect answers in BLUE.}
\end{center}

\vspace{0.75em}

\QBlock{Q1}{Which modeling/approximation statement is the most justified in a typical introductory-physics calculation?}{
  \Wrong{You can model a spinning ceiling fan blade as a point mass when calculating the torque on the fan motor.}{This ignores the blades extended shape and mass distribution, which directly affects rotational inertia and torque. Treating an extended rotating object as a point mass is the wrong model for torque/rotation problems.}
  \Wrong{When analyzing the bending stress inside a diving board, you can treat the board as a point particle located at its center.}{Bending stress depends on how forces and geometry vary along the board, so a point-particle model throws away the essential spatial information. This is a common confusion between translational motion models and deformation/elasticity models.}
  \Wrong{If air drag is described as ``small,'' you must still include it explicitly; you are never allowed to approximate the motion as if in vacuum.}{Neglecting drag is often modeled as motion ``in vacuum'' (or with only gravity), which can be an excellent approximation. The mistake is thinking approximations are never allowed in physics modeling.}
  \Correct{To estimate a cars change in position while it travels $12\,\text{km}$ along a straight highway segment, it is usually reasonable to treat the car as a point mass.}{For large-scale translational motion, the objects size (a few meters) is negligible compared with the distance traveled ($12\,\text{km}$). In that case we track the center-of-mass position and can model the car as a point mass without significantly affecting $\Delta \vec r$ calculations.}
}

\QBlock{Q2}{A student records a small electric charge as $7.5\,\text{nC}$ (nanocoulombs). Which of the following equations is correct?}{
  \Correct{$7.5\,\text{nC} = 7.5\times 10^{-9}\,\text{C}$}{Metric prefixes encode powers of ten: $\text{n} = 10^{-9}$. Therefore $7.5\,\text{nC} = 7.5\times 10^{-9}\,\text{C}$, which is the same physical quantity written in different units.}
  \Wrong{$7.5\,\text{nC} = 7.5\times 10^{9}\,\text{C}$}{This incorrectly flips the prefix; nano means $10^{-9}$, not $10^{+9}$. Its a common sign/exponent error with metric prefixes.}
  \Wrong{$1\,\text{nC} = 10^{-6}\,\text{C}$}{This confuses nano ($10^{-9}$) with micro ($10^{-6}$). The mistake is using the wrong metric prefix.}
  \Wrong{$7.5\,\text{nC} = 0.75\times 10^{-9}\,\text{C}$}{This is off by a factor of 10 due to a decimal-place shift. Its a common scientific-notation slip when rewriting numbers.}
}

\QBlock{Q3}{Which option lists an incorrect SI unit for the stated physical quantity?}{
  \Wrong{Power: watt $(\text{W})$}{Power is correctly measured in watts, with $1\,\text{W}=1\,\text{J/s}$.}
  \Wrong{Pressure: pascal $(\text{Pa}) = \text{N/m}^2$}{Pressure is correctly measured in pascals, and $1\,\text{Pa}=1\,\text{N/m}^2$.}
  \Wrong{Momentum: $\text{kg}\cdot\text{m/s}$}{Momentum $\vec p = m\vec v$ has units $\text{kg}\cdot\text{m/s}$.}
  \Correct{Electric potential (voltage): $\text{C/J}$}{Electric potential is energy per unit charge: $V=\dfrac{U}{q}$. Thus the SI unit is $\text{J/C}$, which is the volt (V), not $\text{C/J}$. Using $\text{C/J}$ is an inversion error (reciprocal unit mistake).}
}

\QBlock{Q4}{A wildlife biologist uses a 2D coordinate grid (units: kilometers) to track a tagged bird. At a certain moment, the birds position relative to the nesting site at the origin is $\vec r = 4.5\,\hat{i} - 7.2\,\hat{j}\,\text{km}$. Which statement is correct?}{
  \Wrong{The bird is $11.7\,\text{km}$ from the nest.}{This adds components ($4.5+7.2$) instead of using the Euclidean distance formula. Thats a common mistake of treating perpendicular components as if they were colinear.}
  \Wrong{The bird is $7.2\,\text{km}$ from the nest.}{This uses only the $y$-component magnitude as the distance, ignoring the $x$-component. The distance depends on both components via the Pythagorean theorem.}
  \Correct{The bird is approximately $8.5\,\text{km}$ from the nest.}{Distance from the origin is the magnitude $|\vec r|=\sqrt{x^2+y^2}=\sqrt{(4.5)^2+(-7.2)^2}\,\text{km}$. This equals $\sqrt{20.25+51.84}=\sqrt{72.09}\approx 8.49\,\text{km}\approx 8.5\,\text{km}$ (verified by calculation).}
  \Wrong{The bird is $4.5\,\text{km}$ from the nest.}{This uses only the $x$-component magnitude as the distance. The error is ignoring the perpendicular $y$-component contribution to $|\vec r|$.}
}

\QBlock{Q5}{A cart moves along the $x$-axis with displacement $\vec s(t)=\big((0.80\,\text{m/s}^4)t^4-(3.6\,\text{m/s}^2)t^2\big)\hat{i}$. Which statement is correct?}{
  \Wrong{This motion has constant acceleration.}{Because $s(t)$ contains a $t^4$ term, the second derivative $a(t)$ depends on time. The misconception is assuming any polynomial motion implies constant acceleration.}
  \Correct{The magnitude of the acceleration at $t=0$ is $7.2\,\text{m/s}^2$.}{Differentiate: $v(t)=\dfrac{ds}{dt}=(3.2t^3-7.2t)\hat{i}$ and $a(t)=\dfrac{dv}{dt}=(9.6t^2-7.2)\hat{i}$. At $t=0$, $a(0)=-7.2\,\text{m/s}^2$, so the magnitude is $|a(0)|=7.2\,\text{m/s}^2$.}
  \Wrong{The magnitude of the acceleration is $3.6\,\text{m/s}^2$.}{The $3.6$ coefficient is in the displacement function, not the acceleration. This is the common error of reading a position coefficient as if it were an acceleration.}
  \Wrong{The acceleration at $t=0$ is $0\,\text{m/s}^2$ because the displacement is zero at $t=0$.}{Acceleration depends on the second derivative of position, not on whether $s(0)=0$. The misconception is conflating position value with derivatives (velocity/acceleration).}
}

\QBlock{Q6}{A warehouse robot moves in a straight line from point $A$ at $\vec r_A=-3.5\hat{i}+12.0\hat{j}\,\text{m}$ to point $B$ at $\vec r_B=4.5\hat{i}+7.0\hat{j}\,\text{m}$ in $4.0\,\text{s}$. Which statement is correct?}{
  \Correct{The distance traveled is approximately $9.4\,\text{m}$.}{Because the path is straight, the distance equals the magnitude of the displacement $\Delta\vec r=\vec r_B-\vec r_A=(8.0\hat{i}-5.0\hat{j})\,\text{m}$. Thus $|\Delta\vec r|=\sqrt{8.0^2+(-5.0)^2}=\sqrt{89}\approx 9.43\,\text{m}\approx 9.4\,\text{m}$ (verified by calculation).}
  \Wrong{The speed at $t=0$ must be $\dfrac{9.4\,\text{m}}{4.0\,\text{s}}\approx 2.35\,\text{m/s}$.}{$\dfrac{\text{distance}}{\text{time}}$ gives the average speed, not necessarily the initial speed. The error is assuming constant speed without it being stated.}
  \Wrong{The acceleration at $t=0$ must be $0\,\text{m/s}^2$ because the motion is along a straight line.}{Straight-line motion can still have nonzero acceleration (speeding up/slowing down). The misconception is confusing ``straight path'' with ``constant velocity.''}
  \Wrong{None of the above is correct.}{Option A follows directly from $|\Delta\vec r|$ for a straight path, so ``none'' is incorrect.}
}

\QBlock{Q7}{A dolphin leaps out of the water with speed $v_0=14.0\,\text{m/s}$ at an angle $\theta=40^\circ$ above the horizontal (ignore air resistance). What is the dolphins speed at maximum height?}{
  \Wrong{$14.0\,\text{m/s}$}{At the top, the vertical component is zero, so the speed is less than $v_0$ unless $\theta=0$. The mistake is assuming speed stays constant in projectile motion.}
  \Wrong{$9.00\,\text{m/s}$}{This equals $v_0\sin\theta$, the initial vertical component, which becomes zero at maximum height. The misconception is swapping sine and cosine roles for components.}
  \Wrong{$0\,\text{m/s}$}{Only the vertical velocity becomes zero at the top; the horizontal component remains. The mistake is thinking the projectile ``stops'' at maximum height.}
  \Correct{$10.7\,\text{m/s}$}{In projectile motion with no drag, the horizontal velocity stays constant: $v_x=v_0\cos\theta$. At maximum height, $v_y=0$, so the speed is $v=|v_x|=v_0\cos\theta=14.0\cos 40^\circ\approx 10.7\,\text{m/s}$ (verified by calculation).}
}

\QBlock{Q8}{A quadcopters initial velocity is $\vec v_0=(18.4\,\text{m/s})\hat{i}$. It then experiences a constant acceleration $\vec a = (-1.6\,\text{m/s}^2)\hat{i}+(3.8\,\text{m/s}^2)\hat{j}$. Which statement is correct?}{
  \Wrong{The $x$-component of the velocity remains constant.}{Since $a_x\neq 0$, the $x$-component changes as $v_x(t)=v_{0x}+a_x t$. The misconception is that only changes in direction affect components.}
  \Correct{Both the direction of motion and the speed change with time.}{Because $a_x\neq 0$, the magnitude of the $x$-velocity changes, which affects the speed. Because $a_y\neq 0$, a $y$-component of velocity develops, so the velocity vectors direction changes; thus both direction and speed change.}
  \Wrong{The speed stays constant, but the direction changes.}{A constant speed would require acceleration always perpendicular to velocity. Here $a_x$ is parallel/antiparallel to the initial velocity, so the speed must change.}
  \Wrong{The $y$-component of the velocity remains zero because the initial $y$-velocity is zero.}{With $a_y\neq 0$, the $y$-velocity becomes $v_y(t)=a_y t$, so it does not remain zero. The misconception is thinking an initially zero component must stay zero.}
}

\QBlock{Q9}{Two point charges exert an electrostatic force on each other. If the magnitude of each charge is multiplied by $4$ and their separation distance is multiplied by $2$, what happens to the magnitude of the force?}{
  \Wrong{It stays the same.}{Coulombs law scales as $F\propto \dfrac{q_1 q_2}{r^2}$; the factors do not cancel here. The error is assuming ``both changed'' means no net change.}
  \Wrong{It decreases by a factor of $4$.}{This treats the distance factor $(2)$ as if it were not squared, or it applies scaling in the wrong direction. The misconception is mishandling the $r^2$ dependence.}
  \Wrong{It increases by a factor of $16$.}{This accounts for multiplying each charge by $4$ (giving $\times 16$) but forgets that doubling the distance reduces force by $2^2=4$.}
  \Wrong{It decreases by a factor of $16$.}{This incorrectly combines both effects as reductions, rather than recognizing that increasing charges increases the force. It reflects a sign/direction confusion in proportional reasoning.}
  \Correct{It increases by a factor of $4$.}{By Coulombs law, $F = k\dfrac{|q_1 q_2|}{r^2}$. Scaling gives $F' = k\dfrac{|(4q_1)(4q_2)|}{(2r)^2}=k\dfrac{16|q_1 q_2|}{4r^2}=4F$, so the force quadruples.}
}

\QBlock{Q10}{For a planet of mass $M$ and radius $R$, the escape speed is $v=\sqrt{\dfrac{2GM}{R}}$. A second planet has mass $4M$ and radius $\dfrac{R}{2}$. By what factor is the escape speed from the second planet larger than from the first?}{
  \Correct{$\sqrt{8}\approx 2.83$}{Use the ratio $\dfrac{v_2}{v_1}=\sqrt{\dfrac{(2G)(4M)/(R/2)}{(2G)M/R}}$. Inside the square root the factor is $\dfrac{4}{1/2}=8$, so $v_2/v_1=\sqrt{8}\approx 2.83$.}
  \Wrong{$2$}{This treats only the mass change ($\sqrt{4}=2$) and ignores the radius change. The misconception is forgetting the $1/\sqrt{R}$ dependence.}
  \Wrong{$\sqrt{2}\approx 1.41$}{This typically comes from incorrectly combining factors, e.g., using $\sqrt{4/(2)}$ instead of $\sqrt{4/(1/2)}$. Its an algebraic inversion error with $R/2$.}
  \Wrong{$8$}{This forgets the square root in $v=\sqrt{2GM/R}$. The misconception is assuming escape speed scales linearly with $GM/R$ instead of with its square root.}
}

\QBlock{Q11}{A constant force $\vec F = (-2.8\,\text{N})\hat{i} + (6.5\,\text{N})\hat{j}$ acts on a particle while it moves along the $y$-axis from $y=-1.2\,\text{m}$ to $y=4.6\,\text{m}$. What is the work done by the force?}{
  \Wrong{$-16.2\,\text{J}$}{This uses the $x$-component $F_x$ even though $\Delta x=0$. The misconception is thinking any force component contributes regardless of displacement direction.}
  \Wrong{$41.0\,\text{J}$}{This likely comes from using the force magnitude $|\vec F|$ times distance instead of the dot product $\vec F\cdot\Delta\vec r$.}
  \Correct{$37.7\,\text{J}$}{Work is $W=\int \vec F\cdot d\vec r$. Since motion is purely along $y$, $d\vec r = dy\,\hat{j}$, so $W=\int F_y\,dy = F_y\Delta y = (6.5)(4.6-(-1.2))\approx 37.7\,\text{J}$ (verified by calculation).}
  \Wrong{$45.0\,\text{J}$}{This corresponds to using an incorrect displacement (wrong $\Delta y$) or rounding too early. The underlying error is not applying $W=\vec F\cdot\Delta\vec r$ carefully.}
}

\QBlock{Q12}{In a crash-test, a $6.2\,\text{kg}$ cart moves toward a barrier at $3.6\,\text{m/s}$ and rebounds straight back at $2.4\,\text{m/s}$. The contact time is $0.15\,\text{s}$. What is the magnitude of the average force on the cart during the collision?}{
  \Wrong{$49.6\,\text{N}$}{This comes from using $|\Delta v|=3.6-2.4=1.2\,\text{m/s}$, incorrectly ignoring the direction reversal. The misconception is failing to treat velocity as a signed quantity.}
  \Correct{$248\,\text{N}$}{Use impulse-momentum: $\vec F_{\text{avg}}\Delta t = m\Delta \vec v$. Taking the initial direction as positive, $v_i=+3.6$ and $v_f=-2.4$, so $\Delta v=v_f-v_i=-6.0\,\text{m/s}$; hence $|F_{\text{avg}}|=\dfrac{m|\Delta v|}{\Delta t}=\dfrac{6.2\times 6.0}{0.15}\approx 2.48\times 10^2\,\text{N}$.}
  \Wrong{$149\,\text{N}$}{This typically results from dividing by the wrong time or using only one of the speeds. The error is an arithmetic/impulse setup mistake rather than applying $m\Delta v$ correctly.}
  \Wrong{$0\,\text{N}$}{A rebound requires a significant momentum change, so the average force cannot be zero. The misconception is confusing ``starts and ends moving'' with ``no net impulse.''}
}

\QBlock{Q13}{A $62\,\text{kg}$ skateboarder is initially at rest on level ground. A friend throws a $2.5\,\text{kg}$ backpack horizontally toward them at $25\,\text{mph}$, and the skateboarder catches it (they move together afterward). What is their speed immediately after the catch? (Use $1\,\text{mph}=0.44704\,\text{m/s}$.)}{
  \Correct{$0.43\,\text{m/s}$}{Treat the catch as a perfectly inelastic collision and conserve momentum: $m_b v_b = (m_s+m_b)v_f$. Convert $25\,\text{mph}$ to $v_b=25(0.44704)\approx 11.18\,\text{m/s}$, so $v_f=\dfrac{(2.5)(11.18)}{62+2.5}\approx 0.433\,\text{m/s}\approx 0.43\,\text{m/s}$ (verified by calculation).}
  \Wrong{$0.50\,\text{m/s}$}{This is a plausible estimate but comes from rounding too aggressively or using an incorrect combined mass. The misconception is treating approximations as exact without checking the conservation equation carefully.}
  \Wrong{$4.3\,\text{m/s}$}{This often comes from forgetting to divide by the total mass $(m_s+m_b)$. It reflects the mistake of conserving speed rather than momentum.}
  \Wrong{$11.2\,\text{m/s}$}{This is essentially the backpacks speed after unit conversion and ignores the skateboarders much larger mass. The misconception is assuming the heavier person doesnt affect the final velocity in an inelastic catch.}
}

\end{document}
